\section{证据与阅读纪律}

本章首先仍然是一个建模章节。它的形式主张只是在下列前提之下被证明：
\begin{enumerate}[nosep]
  \item 本章自己的定义；
  \item 第~\ref{ch:formal-language} 章给出的单调障碍支架；
  \item 本章引入的决策负载假设。
\end{enumerate}

它\emph{并不}主张已经存在实验室层面的证明，说明每一个人类计划都会在每一个领域降低行动成本。本章证明的是模型内部能够推出的内容，仅此而已。

\section{推荐阅读}

\begin{readingbox}
\textbf{基础}
\begin{itemize}[nosep]
  \item Steven Strogatz, \emph{Nonlinear Dynamics and Chaos}，用于理解机制变更的逻辑，以及为什么微小的结构改变可能比单纯加力更重要。
  \item Eugene Izhikevich, \emph{Dynamical Systems in Neuroscience}，用于获得比一般自我管理语言更有纪律的状态转变词汇。
\end{itemize}
\textbf{模型纪律}
\begin{itemize}[nosep]
  \item 在阅读本章之前，先回看第~\ref{ch:formal-language} 章，不要把它误读成已经建立了一条已校准的行为定律。它并没有。
  \item 把第~\ref{ch:dynamics} 章与本章一起读：计划之所以有价值，是因为它在决策窗口开启之前就改变了后续的障碍结构。
\end{itemize}
\textbf{边界}
\begin{itemize}[nosep]
  \item 实证性的意识来源能够说明结构化状态转变值得认真对待，但它们并不是本章计划命题的直接证明。
  \item 为了把计划理解为一种架构控制，并不需要诉诸前沿物理理论。
\end{itemize}
\end{readingbox}

\begin{summarybox}
本章现在少了一些假装，多了一些真正的工作。它定义了计划会改变什么，证明了在模型的单调假设下，解决含糊性和增加准备会降低未来障碍，证明了一个有条件的二元执行优势，并把更广泛的行为主张降级为一种有纪律的解释。这比本章早先的修辞更窄，但离教材式的诚实更近。
\end{summarybox}